\documentclass{jmlr}

\usepackage{booktabs}
\usepackage{amsmath}
\usepackage{graphicx}
\usepackage{xcolor}

\jmlrvolume{--}
\jmlryear{2026}
\jmlrsubmitted{--}
\jmlrpublished{--}

\title{Reproducibility Audit and Regeneration of Table 1 (DeFi Runtime
Comparison) in the HypatiaX Benchmark Report}

\author{\Name{HypatiaX Benchmark Reproducibility Audit}}

\editor{--}

\begin{document}

\jmlrmaketitle

\begin{abstract}
The manuscript's DeFi runtime table (``Table 1'') reports a neural-network
(NN) baseline of 3.0\,s mean / 2.7\,s median, a hybrid arm of 6.8\,s mean /
1.7\,s median, and a resulting 1.73$\times$ / 1.64$\times$ speedup of hybrid
over NN. We regenerate this table mechanically, with no manual number
editing, directly from every raw per-task result file supplied alongside
the manuscript (\texttt{generate\_table1.py}), across six independent
reconstructions of the underlying data pool. In all six, the reported
direction of the effect reverses: the hybrid arm is 4.0$\times$--9.4$\times$
\emph{slower} than the pure-NN arm, never faster, and the absolute NN and
hybrid timings differ from the manuscript by roughly an order of magnitude.
We further document a data-completeness defect in the raw files themselves
(seeds other than 42 contain only 5 of 74 tasks, and the same 5 tasks in
every file) that must be resolved before any pooled multi-seed statistic
can be trusted. On the strength of this evidence we reclassify the
1.73$\times$/1.64$\times$ speedup claim out of the ``no current
contradiction, low-risk footnote'' bucket in which it was previously
grouped (audit item 12) and into a primary, blocking finding. The remaining
items originally grouped under item 12 --- the 89.2\%/62.2\% DeFi accuracy
figures, the Planck blackbody cross-split divergence, PCA
\texttt{random\_state} documentation, and the second \texttt{exp2\_run.log}
--- are not touched by the data examined here and are retained as
lower-risk open follow-up items, listed in a closing footnote.
\end{abstract}

\section{Scope}

This note concerns exactly one artifact: the DeFi wall-clock runtime table
in the manuscript (``Table 1''), and the thirteen raw JSON result files and
one regeneration script supplied for the audit:

\begin{itemize}
\item \texttt{hypatiax\_defi\_benchmark\_v3\_results\_seed\{42,99,123,777,2024\}.json}
  (``aggressive 40/60 split'', v3c benchmark),
\item \texttt{hypatiax\_defi\_benchmark\_pca\_results\_seed\{42,99,123,777,2024\}.json}
  and the unsuffixed \texttt{hypatiax\_defi\_benchmark\_pca\_results.json}
  (``PCA-directed 40/60 split''),
\item \texttt{generate\_table1.py}, a mechanical table-generation script
  whose own header notes that the manuscript's numbers do not reproduce
  from the two seed-42 files it was originally audited against.
\end{itemize}

No hand-editing of any number was performed anywhere in this note; every
statistic below is the direct, unmodified console/LaTeX output of
\texttt{generate\_table1.py} run against the files as supplied.

\section{Data-completeness defect found in the raw files}
\label{sec:completeness}

Before regenerating the table, we checked task counts and task identity
across files. Two things are true of every seed-suffixed file except
seed 42, for \emph{both} benchmark variants:

\begin{enumerate}
\item Each contains only 5 tasks, versus 74 in the seed-42 files
  (6.8\% coverage).
\item All five are the \emph{same} five \texttt{equation\_id} values in
  every file --- \textit{Annualised Portfolio tracking error, Portfolio
  Sharpe Ratio, Portfolio VaR for two correlated, Correlated Portfolio VaR,
  Portfolio Expected Shortfall for correlated} --- regardless of seed or
  benchmark variant.
\end{enumerate}

This is not consistent with five independently-seeded full sweeps that
happened to be interrupted at different points; it is consistent with a
single truncated task ordering that a sweep across all five seeds hit and
stopped at identically, every time. Whatever run produced these files, the
seeds 99/123/777/2024 do not currently constitute evidence about
run-to-run variance across the full 74-task suite --- only about five easy
portfolio-risk tasks. \textbf{This should be treated as a blocking data
gap for the multi-seed sweep recommended in \texttt{generate\_table1.py}'s
own header, not folded into an aggregate statistic as if it were complete
data.}

A second, unrelated hazard: \texttt{hypatiax\_defi\_benchmark\_pca\_
results.json} (no seed suffix) and \texttt{hypatiax\_defi\_benchmark\_pca\_
results\_seed42.json} both report \texttt{seed=42}, but are not the same
data --- the former has all 74 tasks, the latter only the same
five-task truncated subset described above. A future reproduction attempt
that assumes ``the seed-42 file'' is a unique, unambiguous artifact will
silently pick up either a 74-task or a 5-task sample depending on which
file it globs.

\section{Regenerated Table 1, all reconstructions side by side}

Given the ambiguity documented in Section~\ref{sec:completeness}, we did
not pick a single ``correct'' pooling and report only that number.
Instead we ran \texttt{generate\_table1.py} over every reconstruction of
the data that a careful reader of the manuscript's methods section or of
the script's own header could plausibly arrive at: single full seed-42
runs for each variant, the exact multi-seed pooling recipe documented in
the script's header (using the files exactly as named), and a pooling that
substitutes the complete unsuffixed PCA file for its truncated
seed-suffixed counterpart. Table~\ref{tab:timing} reports all six side by
side; Table~\ref{tab:timing_llm_routed} restricts the hybrid arm to
LLM-routed tasks only, matching how the manuscript's hybrid arm is
presumably defined.

\begin{table}[htbp]
\centering
\small
\caption{Wall-clock time per task (seconds), mean / median, regenerated
directly from raw per-task result files with no manual editing. ``Speedup''
compares the neural baseline to hybrid; the manuscript claims
1.73$\times$/1.64$\times$ (hybrid faster). Every reconstruction here shows
the opposite sign.}
\label{tab:timing}
\begin{tabular}{lrrrl}
\toprule
Reconstruction & Pure LLM (s) & Neural MLP (s) & Hybrid, all (s) & Speedup (mean / median) \\
\midrule
v3, seed 42 only (full, $n{=}74$) & 8.94 / 8.01 & 0.36 / 0.36 & 2.03 / 1.56 & 5.59$\times$ / 4.28$\times$ SLOWER \\
PCA, seed 42 only (full file, $n{=}74$) & 8.85 / 7.95 & 0.37 / 0.37 & 1.94 / 1.48 & 5.22$\times$ / 4.00$\times$ SLOWER \\
PCA, seed 42 only (truncated file, $n{=}5$) & 7.68 / 7.19 & 0.27 / 0.27 & 2.35 / 2.28 & 8.68$\times$ / 8.51$\times$ SLOWER \\
v3, multiseed pooled (recipe as documented, $n{=}94$) & 8.88 / 8.01 & 0.36 / 0.36 & 1.98 / 1.53 & 5.52$\times$ / 4.19$\times$ SLOWER \\
PCA, multiseed pooled (recipe as documented, $n{=}25$) & 8.76 / 8.31 & 0.27 / 0.27 & 2.43 / 2.56 & 8.91$\times$ / 9.36$\times$ SLOWER \\
PCA, multiseed pooled (full seed-42 file substituted, $n{=}94$) & 8.89 / 8.08 & 0.35 / 0.37 & 2.05 / 1.60 & 5.84$\times$ / 4.31$\times$ SLOWER \\
\bottomrule
\end{tabular}
\end{table}

\begin{table}[htbp]
\centering
\small
\caption{As Table~\ref{tab:timing}, restricted to tasks where the hybrid
arm's own \texttt{decision} field routed to the LLM rather than the NN.}
\label{tab:timing_llm_routed}
\begin{tabular}{lrrl}
\toprule
Reconstruction ($n$ LLM-routed / total) & Neural MLP (s) & Hybrid, LLM-routed (s) & Speedup (mean / median) \\
\midrule
v3, seed 42 only (69/74) & 0.36 / 0.36 & 1.79 / 1.54 & 4.91$\times$ / 4.23$\times$ SLOWER \\
PCA, seed 42 only, full file (69/74) & 0.37 / 0.37 & 1.67 / 1.44 & 4.50$\times$ / 3.88$\times$ SLOWER \\
PCA, seed 42 only, truncated file (5/5) & 0.27 / 0.27 & 2.35 / 2.28 & 8.68$\times$ / 8.51$\times$ SLOWER \\
v3, multiseed pooled, documented recipe (89/94) & 0.36 / 0.36 & 1.78 / 1.49 & 4.98$\times$ / 4.09$\times$ SLOWER \\
PCA, multiseed pooled, documented recipe (25/25) & 0.27 / 0.27 & 2.43 / 2.56 & 8.91$\times$ / 9.36$\times$ SLOWER \\
PCA, multiseed pooled, full file substituted (89/94) & 0.35 / 0.37 & 1.85 / 1.53 & 5.27$\times$ / 4.12$\times$ SLOWER \\
\bottomrule
\end{tabular}
\end{table}

No timeouts were flagged on the NN arm in any file (\texttt{timed\_out}
is \texttt{false} throughout), so the discrepancy is not an artifact of
excluded failed runs.

\section{Findings}

\begin{enumerate}

\item \textbf{The manuscript's absolute timings do not match any raw file.}
The manuscript reports NN mean/median of 3.0\,s/2.7\,s. Every raw file
audited shows NN times of 0.27--0.39\,s per task, roughly an order of
magnitude faster. This is consistent with the version-drift hypothesis
already recorded in \texttt{generate\_table1.py}'s header: the reported
numbers plausibly came from a heavier, superseded NN implementation, not
the NN training code currently shipped with the benchmark scripts. We find
no defect in the current scripts' NN training code itself that would
explain the gap.

\item \textbf{The manuscript's speedup direction is wrong, robustly.} Across
all six reconstructions in Table~\ref{tab:timing} --- spanning both
benchmark variants, both seed-42 files where two exist, and both the
literal documented pooling recipe and a corrected version of it --- the
hybrid arm is slower than the pure-NN arm by a factor of 4.0$\times$ to
9.4$\times$. It is never faster, in either the ``all tasks'' or
``LLM-routed-only'' framing. This is not sensitive to which file happens
to be used for seed 42, nor to whether the incomplete non-42 seed files are
included.

\item \textbf{The non-42 seed files are not fit for a variance claim.} As
documented in Section~\ref{sec:completeness}, seeds 99/123/777/2024 each
contain the identical 5-task truncated subset, for both benchmark variants.
Any claim of stability or variance of the reported metrics across seeds
that relies on these files should be withdrawn until the sweep is re-run to
completion (74/74 tasks per seed, as the script's own documented recipe
specifies).

\item \textbf{Filename ambiguity is a reproducibility hazard.}
\texttt{hypatiax\_defi\_benchmark\_pca\_results.json} and
\texttt{hypatiax\_defi\_benchmark\_pca\_results\_seed42.json} share a seed
value but not a task set (74 vs.\ 5 tasks). We recommend the re-run
recipe write only seed-suffixed, task-count-stamped filenames going
forward, so no two files can silently disagree about which one is
canonical for a given seed.

\end{enumerate}

\section{Disposition of audit item 12}

Item 12 of the prior audit pass grouped four loose ends together as
``none currently contradicted, all lower-risk to leave as open follow-up
items with a footnote'': the 1.73$\times$ speedup and 89.2\%/62.2\% DeFi
figures, the Planck blackbody cross-split divergence, PCA
\texttt{random\_state} documentation, and a second \texttt{exp2\_run.log}.
That grouping is no longer accurate for one member of the group.

The data examined in this note directly and repeatedly contradicts the
1.73$\times$/1.64$\times$ DeFi speedup claim: it is not merely
unverified, it is reproducibly false under every reconstruction of the raw
data we could construct, including the literal recipe the benchmark
scripts' own documentation recommends. \textbf{We therefore remove the
1.73$\times$/1.64$\times$ speedup claim from item 12 and elevate it to a
primary finding of this note (Table~\ref{tab:timing}, Findings 1--2
above), requiring correction in the manuscript rather than a footnote.}

We did not independently investigate the paired 89.2\%/62.2\% DeFi accuracy
figures beyond what falls out of the same raw files (they are not runtime
figures and are not computed by \texttt{generate\_table1.py}); we make no
claim about them either way and leave them, along with the remaining three
items, in the open-follow-up bucket.\footnote{Open follow-up items
retained from audit item 12, unchanged by this note: (a) the
89.2\%/62.2\% DeFi accuracy figures cited alongside the speedup claim --
not examined here, as they fall outside the runtime files audited; (b) the
Planck blackbody cross-split divergence; (c) documentation of the PCA
benchmark's \texttt{random\_state} handling; (d) the status of a second,
previously referenced \texttt{exp2\_run.log}. None of these four is
contradicted by any evidence in this note; they remain open, lower-risk
items pending a dedicated pass.}

\section{Recommendation}

Replace the manuscript's Table 1 with numbers regenerated by
\texttt{generate\_table1.py} against a completed re-run: all five seeds,
all 74 tasks each, for both the v3c and PCA-directed splits, using the
exact invocation given in the script's own header. Until that re-run
exists, the only currently supportable statement about relative NN/hybrid
runtime is the one this note's data supports: \emph{hybrid is slower than
the pure neural baseline, by roughly 4--9$\times$ depending on pooling,
in every file available.}

\end{document}
